\specialsection{拓扑向量空间}

当 K 不同于 R 或 C 时，我们把 K 上向量空间 F 的半范数（分别地，范数）称为 F 上的超半范数（分别地，作为范数的超半范数）（Esp. Vect. Top.，第 II 章，第 2 版，第 1 节）。若 γ 是 F 上的半范数且 x ∈ F，有时也写作 \textbar\textbar x\textbar\textbar γ 而不写 γ(x)。

我们把拓扑可由单个范数定义的 K 上拓扑向量空间称为可赋范空间，把完备的可赋范空间 \(^{1}\) 称为巴拿赫空间。我们把拓扑可由一族半范数定义的 K 上拓扑向量空间称为多范数空间；当 K = R 或 C 时，这一概念与局部凸空间的概念一致。

若 E 和 F 是多范数空间，我们用 \(\mathcal{L}_{m}(\mathrm{E};\mathrm{F})\) 表示从 \(m\) 到 F 的连续 \(\mathbf{E}^{m}\) 线性映射空间，并赋予其在 \(\mathbf{E}^{m}\) 的有界子集上一致收敛的拓扑；它是一个多范数空间。此外，若 E 可赋范，\(\gamma\) 是 F 上的连续半范数，且 \(u\in\mathcal{L}_{m}(\mathrm{E};\mathrm{F})\)，则以 \(\|u\|_{\gamma}\) 表示满足下列条件的实数 \(a\geqslant0\) 的下确界：

\[
\left\| u (x _ {1}, \dots , x _ {m}) \right\| _ {\gamma} \leqslant a \| x _ {1} \| \dots \| x _ {m} \|
\]

对 E 中所有的 \(x_{1}, \ldots, x_{m}\) 均成立。

